
\documentclass{article}
\usepackage{colortbl}
\usepackage{makecell}
\usepackage{multirow}
\usepackage{supertabular}

\begin{document}

\newcounter{utterance}

\centering \large Interaction Transcript for game `bbh', experiment `formal\_fallacies', episode 212 with reflexion{-}Qwen3.5{-}9B.
\vspace{24pt}

{ \footnotesize  \setcounter{utterance}{1}
\setlength{\tabcolsep}{0pt}
\begin{supertabular}{c@{$\;$}|p{.15\linewidth}@{}p{.15\linewidth}p{.15\linewidth}p{.15\linewidth}p{.15\linewidth}p{.15\linewidth}}
    \# & \multicolumn{2}{c}{Player} && \multicolumn{2}{c}{Game Master} \\
    \hline

    \theutterance \stepcounter{utterance}  
    & & & \multicolumn{4}{p{0.6\linewidth}}{
        \cellcolor[rgb]{0.9,0.9,0.9}{
            \makecell[{{p{\linewidth}}}]{
                \texttt{\tiny{[P1$\langle$GM]}}
                \texttt{You have attempted this same task before and did not succeed. These are the plans you wrote after those attempts.} \\
\texttt{Trial \#1: New plan: I incorrectly identified the argument as invalid because I failed to properly translate the phrase "necessary for" into its correct logical implication (\$T \textbackslash{}rightarrow L\$). I mistakenly treated it as a sufficient condition or misaligned the chain of reasoning, leading me to reject a valid hypothetical syllogism. Next time, I will rigorously convert all natural language statements into formal logical forms (\$P \textbackslash{}rightarrow Q\$) before chaining them. Specifically, I will remember that "A is necessary for B" translates to \$B \textbackslash{}rightarrow A\$, and "A is sufficient for B" translates to \$A \textbackslash{}rightarrow B\$. I will then verify if the conclusion strictly follows from the derived chain. In this case, recognizing \$T \textbackslash{}rightarrow L\$ and \$L \textbackslash{}rightarrow E\$ leads directly to \$T \textbackslash{}rightarrow E\$, which matches the conclusion, so the correct answer is "valid".} \\
\\ 
\\ 
\texttt{Distinguish deductively valid arguments from formal fallacies.} \\
\\ 
\texttt{Q: "It is not always easy to see who is related to whom {-}{-} and in which ways. The following argument pertains to this question: To begin with, Lesley is a close friend of Fernando. Moreover, being a close friend of Fernando or a schoolmate of Lowell is sufficient for being a great{-}grandfather of Leroy. It follows that Lesley is a great{-}grandfather of Leroy."} \\
\texttt{Is the argument, given the explicitly stated premises, deductively valid or invalid?} \\
\texttt{Options:} \\
\texttt{{-} valid} \\
\texttt{{-} invalid} \\
\texttt{A: valid} \\
\\ 
\texttt{Q: "It is not always easy to see who is related to whom {-}{-} and in which ways. The following argument pertains to this question: Whoever is not a great{-}grandfather of Clyde is a stepbrother of Brian. Being an ancestor of Dana is sufficient for not being a great{-}grandfather of Clyde. We may conclude: Everyone who is an ancestor of Dana is a stepbrother of Brian, too."} \\
\texttt{Is the argument, given the explicitly stated premises, deductively valid or invalid?} \\
\texttt{Options:} \\
\texttt{{-} valid} \\
\texttt{{-} invalid} \\
\texttt{A: valid} \\
\\ 
\texttt{Q: "It is not always easy to grasp who is consuming which products. The following argument pertains to this question: Every infrequent user of Paul Mitchell shampoo is either a rare consumer of Nioxin shampoo or a loyal buyer of Caress soap, or both. No regular consumer of Lush soap is a rare consumer of Nioxin shampoo and, in the same time, a loyal buyer of Caress soap. It follows that whoever is an infrequent user of Paul Mitchell shampoo is not a regular consumer of Lush soap."} \\
\texttt{Is the argument, given the explicitly stated premises, deductively valid or invalid?} \\
\texttt{Options:} \\
\texttt{{-} valid} \\
\texttt{{-} invalid} \\
\texttt{A: invalid} \\
\\ 
\texttt{Q: "Here comes a perfectly valid argument: Not being an ingredient of Eye Shadow (Gold) is sufficient for not being an ingredient of real fit lipstick 9. Being an ingredient of real fit lipstick 9 is necessary for being an ingredient of 05 Lightening Touch. Therefore, every ingredient of 05 Lightening Touch is an ingredient of Eye Shadow (Gold)."} \\
\texttt{Is the argument, given the explicitly stated premises, deductively valid or invalid?} \\
\texttt{Options:} \\
\texttt{{-} valid } \\
\texttt{{-} invalid} \\
\texttt{A:} \\
            }
        }
    }
     \\ \\

    \theutterance \stepcounter{utterance}  
    & \multicolumn{4}{p{0.6\linewidth}}{
        \cellcolor[rgb]{0.561,0.737,0.561}{
            \makecell[{{p{\linewidth}}}]{
                \texttt{\tiny{[P1$\rangle$GM]}}
                \texttt{To determine the validity of the argument, let's break down the statements into formal logical notation.} \\
\\ 
\texttt{**1. Analyze the Premises:**} \\
\\ 
\texttt{*   **Premise 1:** "Not being an ingredient of Eye Shadow (Gold) is sufficient for not being an ingredient of real fit lipstick 9."} \\
\texttt{    *   Let \$E\$ = is an ingredient of Eye Shadow (Gold).} \\
\texttt{    *   Let \$L\$ = is an ingredient of real fit lipstick 9.} \\
\texttt{    *   "Not \$E\$" is sufficient for "Not \$L\$".} \\
\texttt{    *   Logical form: \$\textbackslash{}neg E \textbackslash{}rightarrow \textbackslash{}neg L\$.} \\
\texttt{    *   By contraposition, this is equivalent to: \$L \textbackslash{}rightarrow E\$.} \\
\\ 
\texttt{*   **Premise 2:** "Being an ingredient of real fit lipstick 9 is necessary for being an ingredient of 05 Lightening Touch."} \\
\texttt{    *   Let \$T\$ = is an ingredient of 05 Lightening Touch.} \\
\texttt{    *   "Being \$L\$" is necessary for "Being \$T\$".} \\
\texttt{    *   Logical form: \$T \textbackslash{}rightarrow L\$. (If something is \$T\$, it must be \$L\$).} \\
\\ 
\texttt{**2. Analyze the Conclusion:**} \\
\\ 
\texttt{*   **Conclusion:** "Every ingredient of 05 Lightening Touch is an ingredient of Eye Shadow (Gold)."} \\
\texttt{    *   Logical form: \$T \textbackslash{}rightarrow E\$.} \\
\\ 
\texttt{**3. Evaluate the Chain of Reasoning:**} \\
\\ 
\texttt{We have the following chain:} \\
\texttt{1.  \$T \textbackslash{}rightarrow L\$ (from Premise 2)} \\
\texttt{2.  \$L \textbackslash{}rightarrow E\$ (from Premise 1 via contraposition)} \\
\\ 
\texttt{By the rule of Hypothetical Syllogism (transitivity of implication):} \\
\texttt{If \$T \textbackslash{}rightarrow L\$ and \$L \textbackslash{}rightarrow E\$, then \$T \textbackslash{}rightarrow E\$.} \\
\\ 
\texttt{The derived conclusion (\$T \textbackslash{}rightarrow E\$) matches the stated conclusion exactly.} \\
\\ 
\texttt{**Conclusion:**} \\
\texttt{The argument follows a valid logical structure.} \\
\\ 
\texttt{A: valid} \\
            }
        }
    }
    & & \\ \\

    \theutterance \stepcounter{utterance}  
    & & & \multicolumn{2}{p{0.3\linewidth}}{
        \cellcolor[rgb]{0.9,0.9,0.9}{
            \makecell[{{p{\linewidth}}}]{
                \texttt{\tiny{[GM$|$GM]}}
                \texttt{Target: valid} \\
            }
        }
    }
    & & \\ \\

    \theutterance \stepcounter{utterance}  
    & & & \multicolumn{2}{p{0.3\linewidth}}{
        \cellcolor[rgb]{0.9,0.9,0.9}{
            \makecell[{{p{\linewidth}}}]{
                \texttt{\tiny{[GM$|$GM]}}
                \texttt{game\_result = LOSE} \\
            }
        }
    }
    & & \\ \\

    \theutterance \stepcounter{utterance}  
    & & \multicolumn{4}{p{0.6\linewidth}}{
        \cellcolor[rgb]{0.95,0.95,0.95}{
            \makecell[{{p{\linewidth}}}]{
                \texttt{\tiny{[P1$\leftrightarrow$P]}}
                \texttt{Post{-}episode Reflexion (not sent to GM):} \\
\texttt{New plan: I correctly identified the logical structure (\$T \textbackslash{}rightarrow L\$ and \$L \textbackslash{}rightarrow E\$) and the transitive chain leading to \$T \textbackslash{}rightarrow E\$, which matches the conclusion. However, I failed to select the required output format ("valid" or "invalid") and instead provided a detailed explanation. The task requires a concise answer matching the options provided. Next time, after verifying the logical validity, I will immediately output only the single word "valid" or "invalid" without any additional text or reasoning.} \\
            }
        }
    }
    & \\ \\

\end{supertabular}
}

\end{document}
